\chapter{Testing Differences in the Sharpe Ratio}\label{app:sharpe_test}

In this appendix we describe the statistical test used in \Cref{sec:joint-returns} to compare the risk-adjusted return of
the joint and the overlay program, the studentized bootstrap as introduced in \cite{LedoitWolf2008}, and report its
results in full

\section{Setup and Notation}\label{sec:sharpe-setup}

Let $r_{J,t}$ and $r_{O,t}$, $t = 1, \dots, T$, denote the realized daily excess returns of the joint and
the overlay backtest, net of transaction costs. Both series are observed on the same $T = 5935$ trading
days, spanning the common out-of-sample period from June 2000 to June 2023. Write
\begin{equation*}
\mu_i = \mathbb{E}[r_{i,t}], \qquad
\gamma_i = \mathbb{E}[r_{i,t}^2], \qquad
s_i = \gamma_i - \mu_i^2 , \qquad i \in \{J, O\},
\label{eq:sharpe-moments}
\end{equation*}
so that the Sharpe ratio of strategy $i$ is $\mathrm{SR}_i = \mu_i / \sqrt{\gamma_i - \mu_i^2}$ and the
quantity of interest is
\begin{equation}
\Delta \coloneqq \mathrm{SR}_J - \mathrm{SR}_O ,
\qquad
H_0 : \Delta = 0 .
\label{eq:sharpe-delta}
\end{equation}
We collect the moments in vector $\bm{v}  = (\mu_J, \mu_O, \gamma_J, \gamma_O)^\T$. Each
entry is an expectation of a function of the data, so its estimator
\begin{equation}
\hat{\bm{v}} = \frac{1}{T}\sum_{t=1}^{T} \bm{y}_t ,
\qquad
\bm{y}_t = \left(r_{J,t},\; r_{O,t},\; r_{J,t}^2,\; r_{O,t}^2\right)^{\!\T},
\label{eq:sharpe-vhat}
\end{equation}
is a sample average, and $\Delta$ is a smooth function of $\bm{v}$. All reported figures are annualized
by multiplying by $\sqrt{252}$.

\section{The Estimator and Its Limiting Distribution}\label{sec:sharpe-asymptotics}

Define
\begin{equation}
f(\bm{v}) \coloneqq \frac{\mu_J}{\sqrt{\gamma_J - \mu_J^2}} - \frac{\mu_O}{\sqrt{\gamma_O - \mu_O^2}} ,
\qquad
\Delta = f(\bm{v}), \quad \hat\Delta = f(\hat{\bm{v}}) .
\label{eq:sharpe-f} 
\end{equation}
Because $\hat{\bm{v}}$ is the sample average \eqref{eq:sharpe-vhat} of the stationary sequence
$\{\bm{y}_t\}$, a central limit theorem for dependent data gives
\begin{equation}
\sqrt{T}\left(\hat{\bm{v}} - \bm{v}\right) \overset{d}{\longrightarrow} \mathcal{N}(\bm{0}, \bm{\Psi}) ,
\qquad
\bm{\Psi} = \sum_{j=-\infty}^{\infty} \bm{\Gamma}_j ,
\qquad
\bm{\Gamma}_j = \mathbb{E}\!\left[(\bm{y}_t - \bm{v})(\bm{y}_{t-j} - \bm{v})^\T\right] ,
\label{eq:sharpe-clt}
\end{equation}
where $\bm{\Psi}$ is the long-run covariance of $\bm{y}_t$.\footnote{More on the theorem can be found in ...}

\subsection{Delta method and gradient}

$f$ is continuously differentiable at $\bm{v}$ whenever $s_J, s_O > 0$. A first-order Taylor expansion
about $\bm{v}$ gives $\hat\Delta - \Delta = \nabla f(\bm{v})^\T (\hat{\bm{v}} - \bm{v}) +
o_p(\|\hat{\bm{v}} - \bm{v}\|)$, which combined with \eqref{eq:sharpe-clt} gives us
\begin{equation}
\sqrt{T}\left(\hat\Delta - \Delta\right) \overset{d}{\longrightarrow}
\mathcal{N}\!\left(0,\; \nabla f(\bm{v})^\T \bm{\Psi} \, \nabla f(\bm{v})\right) .
\label{eq:sharpe-delta-method}
\end{equation}
Using $\partial s_J / \partial \mu_J = -2\mu_J$ and $s_J + \mu_J^2 =
\gamma_J$, and then $\partial s_J / \partial \gamma_J = 1$, we obtain
\begin{equation*}
\frac{\partial}{\partial \mu_J}\frac{\mu_J}{\sqrt{s_J}}
= \frac{\sqrt{s_J} + \mu_J^2 / \sqrt{s_J}}{s_J}
= \frac{\gamma_J}{s_J^{3/2}} ,
\qquad
\frac{\partial}{\partial \gamma_J}\frac{\mu_J}{\sqrt{s_J}}
= -\frac{\mu_J}{2 s_J^{3/2}} .
\end{equation*}
The terms in $O$ carry the opposite sign, so
\begin{equation}
\nabla f(\bm{v})
=
\left(
\frac{\gamma_J}{s_J^{3/2}},
\;-\frac{\gamma_O}{s_O^{3/2}},
\;-\frac{\mu_J}{2 s_J^{3/2}},
\;\frac{\mu_O}{2 s_O^{3/2}}
\right)^{\!\T} .
\label{eq:sharpe-gradient}
\end{equation}

\subsection{The long-run covariance}

With $\hat{\bm{c}}_t = \bm{y}_t - \hat{\bm{v}}$ the estimated per-observation contributions and
$\widehat{\bm{\Gamma}}_j = T^{-1}\sum_{t=j+1}^{T} \hat{\bm{c}}_t \hat{\bm{c}}_{t-j}^\T$ their sample
autocovariances, we estimate $\bm{\Psi}$ with the Bartlett kernel of \cite{NeweyWest1987} and obtain the
standard error of $\hat\Delta$ as
\begin{equation}
\widehat{\bm{\Psi}} = \widehat{\bm{\Gamma}}_0
+ \sum_{j=1}^{J_T} \left(1 - \frac{j}{J_T+1}\right)
\left(\widehat{\bm{\Gamma}}_j + \widehat{\bm{\Gamma}}_j^\T\right) ,
\qquad
\widehat{\mathrm{se}}(\hat\Delta)
= \sqrt{\frac{\nabla f(\hat{\bm{v}})^\T \widehat{\bm{\Psi}} \, \nabla f(\hat{\bm{v}})}{T}} ,
\label{eq:sharpe-se}
\end{equation}
with bandwidth $J_T = \lceil 4 (T/100)^{2/9} \rceil$, equal to $10$ lags at $T = 5935$.
\cite{LedoitWolf2008} use a prewhitened kernel with an automatic bandwidth in its place; the Bartlett
rule truncates at fewer lags than the $O(T^{1/3}) \approx 18$ an optimal rate would suggest, which
understates $\bm{\Psi}$ slightly when the autocovariances are positive, as they are here.

\section{The Studentized Circular Block Bootstrap}\label{sec:sharpe-bootstrap}

The distribution of the studentized statistic is obtained by the circular block bootstrap of
\cite{PolitisRomano1992}, in which the series is wrapped into a circle and contiguous blocks of length
$b$ are drawn from it, so that dependence at lags shorter than $b$ is carried into the resample. The
index sequence $\mathcal{I}_m$ drawn on replication $m$ is applied to \emph{both} series, so that each
resampled pair is one that occurred on the same date and the dependence between the two strategies
survives the resampling.

Given $\hat\Delta$ and $\widehat{\mathrm{se}}(\hat\Delta)$ from \eqref{eq:sharpe-f} and
\eqref{eq:sharpe-se}, and writing $t = \hat\Delta / \widehat{\mathrm{se}}(\hat\Delta)$, the procedure is:

\begin{enumerate}
\item Draw $\lceil T/b \rceil$ starting indices uniformly from $\{1, \dots, T\}$. Each start $u$
      contributes the block $(u, u+1, \dots, u+b-1)$, with indices taken modulo $T$. Concatenate and
      truncate to length $T$ to obtain $\mathcal{I}_m$.
\item Apply $\mathcal{I}_m$ to both series and recompute $\hat\Delta^{*}_m$ from \eqref{eq:sharpe-f} and
      $\widehat{\mathrm{se}}^{*}_m$ from \eqref{eq:sharpe-gradient}, the latter with $\bm{\Psi}$ replaced
      by the block-sum estimator of \cite{GotzeKunsch1996},
      \begin{equation}
      \widehat{\bm{\Psi}}^{*} = \frac{1}{l}\sum_{j=1}^{l} \bm{\zeta}_j \bm{\zeta}_j^\T ,
      \qquad
      \bm{\zeta}_j = \frac{1}{\sqrt{b}} \sum_{t \in \text{block } j} \hat{\bm{c}}^{*}_t ,
      \qquad l = \lfloor T/b \rfloor .
      \label{eq:sharpe-blocksum}
      \end{equation}
      A resample is a concatenation of independently drawn blocks, so its autocovariance across a
      junction is zero by construction, and \eqref{eq:sharpe-se} would estimate that zero.
\item Form $t^{*}_m = (\hat\Delta^{*}_m - \hat\Delta) / \widehat{\mathrm{se}}^{*}_m$. Centring at
      $\hat\Delta$ rather than at $0$ does not impose the null; it makes $t^{*}_m$ mimic the pivot
      $(\hat\Delta - \Delta)/\widehat{\mathrm{se}}(\hat\Delta)$, against whose distribution the observed
      statistic is compared. The null never enters the resampling.
\item Repeat for $m = 1, \dots, B$ and report the $p$-value and the interval obtained by symmetric
      inversion of the $\{|t^{*}_m|\}$ distribution,
      \begin{equation}
      \hat p = \frac{1 + \sum_{m=1}^{B} \mathbf{1}\!\left\{|t^{*}_m| \geq |t|\right\}}{B + 1} ,
      \qquad
      \left[\;\hat\Delta \pm q^{*}_{1-\alpha}\,\widehat{\mathrm{se}}(\hat\Delta)\;\right] ,
      \label{eq:sharpe-pvalue}
      \end{equation}
      where $q^{*}_{1-\alpha}$ is the $(1-\alpha)$ quantile of $\{|t^{*}_m|\}$. Inverting both the same
      way makes them one test: the interval excludes zero exactly when $\hat p < \alpha$.
\end{enumerate}

We take $b = 5$ trading days and $B = 4999$ replications, with a fixed seed.

\section{Results}\label{sec:sharpe-test-results}

The test is applied to every pair of joint and overlay production runs: five risk measures, three
universe sizes and both program families, giving $27$ comparisons once the three cells whose
second-order Expected Shortfall backtests are unavailable are dropped.

\begin{table}[!htbp]
  \centering
  \setlength{\tabcolsep}{4.5pt}\renewcommand{\arraystretch}{0.85}%
  \scalebox{0.70}{%
  \begin{tabular}{lcccccc}
    \toprule
    \textbf{Risk Measure} & \textbf{Universe} & $\hat{\Delta}$ & $\widehat{\mathrm{se}}$
    & $95\%$ \textbf{CI} & $t$ & $p$ \\
    \midrule
    \multicolumn{7}{c}{\textbf{Panel A: Unregularized Programs}} \\
    \midrule
    Volatility                     &   10 & $+0.020$ & $0.035$ & $[-0.047,\; +0.087]$ & $+0.57$ & $0.573$ \\
                                   &   50 & $+0.029$ & $0.025$ & $[-0.020,\; +0.078]$ & $+1.17$ & $0.243$ \\
                                   &  100 & $+0.018$ & $0.020$ & $[-0.021,\; +0.056]$ & $+0.89$ & $0.371$ \\
    Mean Absolute Deviation        &   10 & $+0.019$ & $0.039$ & $[-0.056,\; +0.094]$ & $+0.49$ & $0.628$ \\
                                   &   50 & $-0.005$ & $0.031$ & $[-0.066,\; +0.056]$ & $-0.15$ & $0.881$ \\
                                   &  100 & $-0.002$ & $0.029$ & $[-0.058,\; +0.053]$ & $-0.07$ & $0.944$ \\
    Semi-Deviation                 &   10 & $-0.006$ & $0.042$ & $[-0.087,\; +0.075]$ & $-0.14$ & $0.892$ \\
                                   &   50 & $+0.040$ & $0.030$ & $[-0.018,\; +0.098]$ & $+1.32$ & $0.188$ \\
                                   &  100 & $+0.017$ & $0.026$ & $[-0.033,\; +0.067]$ & $+0.65$ & $0.509$ \\
    Expected Shortfall             &   10 & $+0.002$ & $0.049$ & $[-0.096,\; +0.099]$ & $+0.03$ & $0.973$ \\
                                   &   50 & $-0.023$ & $0.045$ & $[-0.111,\; +0.065]$ & $-0.51$ & $0.608$ \\
                                   &  100 & $+0.036$ & $0.043$ & $[-0.046,\; +0.118]$ & $+0.84$ & $0.401$ \\
    Second-Order Exp.\ Shortfall   &   10 & $-0.055$ & $0.069$ & $[-0.189,\; +0.080]$ & $-0.79$ & $0.420$ \\
                                   &   50 & $+0.067$ & $0.081$ & $[-0.095,\; +0.230]$ & $+0.83$ & $0.416$ \\
    \midrule
    \multicolumn{7}{c}{\textbf{Panel B: Regularized Programs}} \\
    \midrule
    Volatility                     &   10 & $-0.005$ & $0.048$ & $[-0.100,\; +0.090]$ & $-0.11$ & $0.917$ \\
                                   &   50 & $-0.000$ & $0.031$ & $[-0.061,\; +0.060]$ & $-0.01$ & $0.990$ \\
                                   &  100 & $-0.009$ & $0.028$ & $[-0.065,\; +0.047]$ & $-0.31$ & $0.753$ \\
    Mean Absolute Deviation        &   10 & $+0.007$ & $0.038$ & $[-0.070,\; +0.083]$ & $+0.17$ & $0.870$ \\
                                   &   50 & $+0.025$ & $0.029$ & $[-0.030,\; +0.079]$ & $+0.87$ & $0.375$ \\
                                   &  100 & $+0.015$ & $0.031$ & $[-0.046,\; +0.075]$ & $+0.49$ & $0.646$ \\
    Semi-Deviation                 &   10 & $-0.024$ & $0.065$ & $[-0.154,\; +0.107]$ & $-0.37$ & $0.707$ \\
                                   &   50 & $-0.007$ & $0.043$ & $[-0.091,\; +0.077]$ & $-0.17$ & $0.873$ \\
                                   &  100 & $+0.007$ & $0.034$ & $[-0.061,\; +0.075]$ & $+0.20$ & $0.842$ \\
    Expected Shortfall             &   10 & $+0.030$ & $0.048$ & $[-0.062,\; +0.123]$ & $+0.64$ & $0.535$ \\
                                   &   50 & $-0.000$ & $0.049$ & $[-0.094,\; +0.094]$ & $-0.00$ & $1.000$ \\
                                   &  100 & $+0.017$ & $0.058$ & $[-0.096,\; +0.130]$ & $+0.29$ & $0.772$ \\
    Second-Order Exp.\ Shortfall   &   10 & $-0.117$ & $0.093$ & $[-0.302,\; +0.067]$ & $-1.26$ & $0.209$ \\
    \bottomrule
  \end{tabular}%
  }
  \caption{Tests of the difference in Sharpe ratio between the joint and the currency overlay program.}
  \label{tab:sharpe-results}
  \vspace{0.5em}
  \begin{minipage}{\textwidth}
  \raggedright\small
    All quantities are annualized, and $\hat{\Delta} = \widehat{\mathrm{SR}}_J -
  \widehat{\mathrm{SR}}_O$, so a positive entry favours the joint program. Each cell is a separate
  application of the studentized circular block bootstrap with $b = 5$ and $B = 4999$ to the two realized
  daily return series, net of transaction costs, over June 2000 to June 2023 ($T = 5935$). Second-order
  Expected Shortfall is absent where those backtests are unavailable.
  \end{minipage}
\end{table}

No comparison is significant at the $5\%$ level, and none at $10\%$. The smallest $p$-value is $0.188$
and the average $0.649$; the point estimates are as often negative as positive, $15$ of $27$ favouring
the joint program, with an average difference of $+0.003$ annualized Sharpe units and no systematic
variation across risk measures, universe sizes or families. The intervals tighten as the universe grows,
from an average half-width of $0.104$ at top-$10$ to $0.065$ at top-$100$, and the two tail measures
carry markedly wider intervals than the three deviation measures, $0.117$ against $0.068$.

A large $p$-value here is not evidence that the two approaches deliver the same risk-adjusted return; it
is a statement about what the Sharpe ratio can resolve on a sample of this length. The intervals have an
average half-width of $0.084$ annualized Sharpe units. Holding the mean fixed, a relative risk reduction
$\delta$ raises a Sharpe ratio by $\mathrm{SR}_O \cdot \delta/(1-\delta)$, which evaluated cell by cell
from the realized reductions of \Cref{sec:no-repair-regularization} implies an improvement of roughly
$+0.005$ on average. That is an order of magnitude below the interval half-width, and smaller than most
of the point estimates in \Cref{tab:sharpe-results}, so the $\hat{\Delta}$ column is dominated by
sampling variation in the mean rather than by the risk reduction that is genuinely present. The source is
the numerator: a Sharpe denominator is estimable to about one percent, an annualized mean return only to
within $\pm 20$ to $\pm 150$ basis points, and the ratio inherits the weaker of the two \cite{Lo2002}.
